% 中文摘要

随着5G通信与流媒体技术的快速发展，视频数据已成为全球互联网流量的主体。在有限带宽下如何获得高质量的视觉体验成为亟待解决的关键问题。然而，传统视频编码标准面临着压缩效率与计算复杂度之间的矛盾，端到端深度学习编码方案虽性能优异但难以在边缘设备上部署。

针对上述问题，本文提出了一种基于强化学习的视频编码协处理方案。该方案以NVIDIA Video Encoder (NVENC)硬件编码器为核心，通过引入轻量级码率自适应前处理网络(RAPNet)实现编码前的智能视频增强。系统基于Actor-Critic架构，采用残差学习策略将动作空间约束为对原始帧的微调修正，并设计了基于相对率失真增益的复合奖励函数，在感知质量与传输码率之间实现有效博弈。此外，提出了包含加权MSE损失、非边缘高频抑制损失和全变分平滑损失的频域感知复合损失策略，有效抑制了增强过程中的高频伪影。

实验在1080p分辨率、YUV420格式的5个HEVC标准测试序列上进行了全面验证。结果表明，本文方案在模型参数量仅为139K、单帧推理时间仅5.4ms的前提下，取得了平均$-9.48\%$的BD-LPIPS增益，其中ParkScene和BasketballDrive序列分别达到$-13.45\%$和$-12.84\%$。前处理开销仅占端到端编码时延的约$12\%$，验证了轻量级AI协处理在工业级实时视频编码场景中的可行性。

% 注：关键词已在 my_paper.tex 的 \whusetup 中设置
